% ============================================================
%  sections/data.tex  –  Data, Descriptive Statistics, and Empirical Strategy
% ============================================================

\subsection{Data Source}

The analysis presented in this paper draws on a crowdsourced dataset of marijuana purchases in Mexico assembled by \citet{bejaranoromero2020dataset}.  An anonymous online questionnaire, administered via Qualtrics from February~19 to September~15, 2019 and mainly promoted through the news outlet \emph{Vice en Espa\~{n}ol}, collected self-reported information on individual transactions.  For each purchase, the respondent recorded the quantity acquired, the amount paid, perceived product quality, the state and municipality of purchase, and the purchase date.  The questionnaire also asked about self-reported socioeconomic status, age of first cannabis consumption, purchase frequency, usual purchase amount, the respondent's relationship with the dealer, and basic demographic characteristics.  The raw dataset contains 4{,}775~records.

We supplement the survey with two external spatial and enforcement data sources.  Municipality and state identifiers are linked to INEGI municipal and state polygons from the Marco Geoestad\'{i}stico Integrado; we calculate municipal centroids for purchase and seizure municipalities and full-state centroids for the production-hub states used in the OSM driving-time instrument \citep{inegi2022mgi}.  Road-network driving times are computed with the Open Source Routing Machine (OSRM) using OpenStreetMap (OSM) data \citep{openstreetmap2026,luxen2011osrm}.  Because the OSRM matrices are built from an OSM road-network snapshot queried in April~2026, these variables should be interpreted as persistent road-network connectivity rather than contemporaneous 2019 traffic conditions.  Enforcement intensity is measured with monthly municipality-level data published by M\'{e}xico Unido contra la Delincuencia (MUCD), a Mexican non-governmental organization.  These records report marijuana seizures, crop destruction, and seed confiscation; the seizure-gravity instrument uses only marijuana seizures from the one-month lag window relevant for the estimation sample, December~2018 through August~2019 \citep{mucd2019seizures}.  MUCD's platform compiles official responses obtained through public-information requests to the transparency units of the Secretariat of the Navy (SEMAR), the Secretariat of National Defense (SEDENA), the National Guard (GN), the former Federal Police (PF), and the Attorney General's Office (FGR); MUCD then transcribed, systematized, cleaned, and fused the heterogeneous source files into a municipality-month database \citep{mucd2019seizures}.  These external sources are used only to construct the instrumental-variables diagnostics; the baseline OLS estimates rely on the cleaned transaction-level survey sample.  Appendix~\ref{app:external-data} reports descriptive diagnostics for the MUCD marijuana-seizure records used in the IV diagnostics and the OSM/OSRM driving-time matrices.

These auxiliary sources characterize geography and enforcement, but the core empirical advantage comes from the survey's transaction-level design.  Conventional data sources usually miss the relevant price--quantity margin: administrative price series do not exist for illegal drugs, and household surveys such as ENCODAT measure consumption prevalence but generally do not record the price and quantity of specific purchases \citep{encodat2025}.  Here, respondents report completed transactions, so the data reflect the terms buyers actually faced rather than hypothetical willingness to pay.  That distinction matters for demand estimation: search costs, legal risk, and imperfect information about quality can all affect realized purchases, and \citet{nisbet1972ucla} show that elasticities based on actual purchases can differ substantially from hypothetical-response estimates.

\subsection{Sample Construction}

The analytic sample is built around the variables needed to observe a completed purchase and estimate a revealed-demand elasticity.  Quantity purchased is the outcome, while transaction value and quantity jointly define the unit price faced by the buyer.  Perceived quality is retained because seeded and seedless marijuana are not the same product, and failing to condition on quality would confound price variation with product differentiation.  Purchase date and municipality/state identifiers are required to align transactions with local market conditions, fixed effects, and the external driving-time and enforcement instruments.  Finally, age, gender, and education are included as baseline demographic controls because they are predetermined respondent characteristics that can shape purchasing behavior, income constraints, risk exposure, and access to suppliers.

The survey also records self-reported socioeconomic status, age of first consumption, purchase frequency, usual purchase amounts, and the respondent's relationship with the dealer.  These variables are informative about consumer experience, market relationships, and socioeconomic position, but we do not use them as baseline controls or sample-defining fields.  Frequency, usual amount, and buyer--seller relationship are closer to behavioral intensity or matching outcomes than to predetermined covariates, and conditioning on them could absorb part of the demand response or the negotiation channel through which prices and quantities are jointly determined.  Socioeconomic status is self-classified and coarser than education, so we treat it as contextual rather than as a main control.  Requiring complete responses for these auxiliary items would also reduce the sample without being necessary for the core price--quantity estimand.

We apply a sequence of filters and transformations to arrive at the analytic sample. First, we require nonmissing values in the fields needed by the main specification: state, municipality, age, gender, education, purchase date, perceived quality, quantity, and transaction value. This complete-case requirement removes 1{,}291 of the 4{,}775 raw records. Second, we exclude 7 additional records with implausible age, transaction-value, or quantity entries. Third, although the questionnaire was fielded from February to September~2019, some respondents reported earlier purchases; restricting to transactions dated on or after January~1, 2019 removes 184 recalled purchases. We then construct the unit price per gram as the ratio of transaction value to quantity purchased, $p_i = \text{value}_i / \text{quantity}_i$, and code perceived quality as a binary indicator, with \emph{good} denoting seedless marijuana and \emph{bad} denoting seeded marijuana. After all filters and transformations, the analytic sample comprises 3{,}293 observations, with retained purchases dated January~1--September~11, 2019.

\subsection{Descriptive Statistics}

Table~\ref{tab:descriptive} presents summary statistics for the full sample and by quality tier.  The median transaction involves 14~grams purchased for 250~MXN, yielding a median price of about 21~MXN per gram.  The distributions of both price and quantity are right-skewed, with means substantially exceeding medians and large standard deviations relative to the mean, motivating the log transformation used in the regressions.  The median respondent is~25 years old.

\begin{table}[H]
\centering
\caption{Descriptive statistics}
\label{tab:descriptive}
\begin{threeparttable}
\begin{tabular}{lrrrrr}
\toprule
 & Mean & Median & SD & Min & Max \\
\midrule
\multicolumn{6}{l}{\textit{Panel A: Full sample (N = 3,293)}} \\[3pt]
Price (MXN per gram)   & 41.97 & 21.43 & 88.35 & 0.01 & 3,000.00 \\
Quantity (grams)       & 31.27 & 14.00 & 83.40 & 1.00 & 2,000.00 \\
Transaction value (MXN)& 568.03 & 250.00 & 1,892.43 & 1.00 & 65,000.00 \\
Age (years)            & 26.20 & 25.00 & 6.74 & 18 & 68 \\
\\[6pt]
\multicolumn{6}{l}{\textit{Panel B: Good quality (seedless) (N = 2,636)}} \\[3pt]
Price (MXN per gram)   & 48.39 & 28.57 & 97.12 & 0.01 & 3,000.00 \\
Quantity (grams)       & 31.31 & 14.00 & 83.54 & 1.00 & 2,000.00 \\
Transaction value (MXN)& 663.33 & 320.00 & 2,100.50 & 1.00 & 65,000.00 \\
Age (years)            & 26.38 & 25.00 & 6.75 & 18 & 65 \\
\\[6pt]
\multicolumn{6}{l}{\textit{Panel C: Bad quality (seeded) (N = 657)}} \\[3pt]
Price (MXN per gram)   & 16.24 & 10.00 & 21.29 & 0.25 & 200.00 \\
Quantity (grams)       & 31.13 & 10.00 & 82.90 & 1.00 & 1,000.00 \\
Transaction value (MXN)& 185.68 & 100.00 & 258.32 & 10.00 & 2,500.00 \\
Age (years)            & 25.47 & 24.00 & 6.68 & 18 & 68 \\
\bottomrule
\end{tabular}
\begin{tablenotes}[flushleft]\small
\item \textit{Notes:} Price per gram is computed as transaction value divided by quantity purchased.  Quality classification is self-reported: ``good'' denotes seedless marijuana and ``bad'' denotes seeded marijuana.
\end{tablenotes}
\end{threeparttable}
\end{table}


Quality segmentation is pronounced.  Seedless (\emph{good}) marijuana, which accounts for 80\% of transactions ($N = 2{,}636$), commands a median price nearly three times that of seeded (\emph{bad}) marijuana (28.57 versus 10.00~MXN per gram).  The price distribution for bad-quality product is considerably more compressed (SD~=~21 versus~97), consistent with a less differentiated, commodity-like segment.  Median quantities are also higher for good-quality purchases (14 versus 10~grams), though the means are similar across tiers.

The sample is predominantly male (84\%), with women comprising 14\% and other gender identities 2\%.  Educational attainment skews toward higher levels: 51\% report a bachelor's degree, 37\% high school, 9\% a graduate degree, and only 3\% middle school or below.  These demographics are consistent with the online, media-promoted recruitment channel and should be borne in mind when considering external validity.

Geographically, the data cover all 32~states and 351~municipalities.  Mexico City dominates with 1{,}011~observations (31\%), followed by the State of Mexico (283), Jalisco (226), and Nuevo Le\'{o}n (150).  The temporal distribution peaks in February and March~2019 (1{,}089 and 903~observations, respectively), coinciding with the survey's launch and initial media push, and tapers in subsequent months.

Quantity heaping is evident: 42\% of reported quantities are multiples of~5, 23\% are multiples of~10, and 14\% are close to~28~grams (approximately one ounce).  This pattern, common in self-reported survey data, motivates a robustness check that excludes round-number quantities.

\subsection{Empirical Strategy}

\paragraph{Baseline specification.}
We estimate the price elasticity of demand using a log--log specification:
\begin{equation}\label{eq:loglog}
  \ln Q_i \;=\; \alpha \;+\; \beta \, \ln P_i \;+\; \gamma \, \textit{Good}_i \;+\; \mathbf{X}_i' \boldsymbol{\delta} \;+\; \mu_s \;+\; \varepsilon_i,
\end{equation}
where $Q_i$ is the quantity purchased (grams), $P_i$ is the unit price (MXN per gram), $\textit{Good}_i$ is an indicator for seedless quality, $\mathbf{X}_i$ is a vector of demographic controls (age, gender, and education), $\mu_s$ denotes fixed effects, and $\varepsilon_i$ is the error term.  The coefficient~$\beta$ is the price elasticity of demand: a one-percent increase in price is associated with a $\beta$-percent change in quantity demanded.

We progressively enrich the fixed-effects structure: (i)~no fixed effects, (ii)~state fixed effects ($\mu_s$), (iii)~state plus month fixed effects, (iv)~state-by-month interactions, and (v)~municipality fixed effects.  The motivation for this progression is to absorb increasingly fine-grained unobserved heterogeneity in local market conditions, enforcement intensity, and demand-side factors.  For the pooled specifications we report heteroskedasticity-robust standard errors; once geographic fixed effects are introduced, standard errors are clustered at the state level.  Because there are 32 state clusters, we also report wild cluster bootstrap inference for the benchmark specification.

\paragraph{Quality heterogeneity.}
To allow the demand response to differ across product tiers, we estimate separate regressions for good and bad quality subsamples, and a pooled specification that interacts $\ln P_i$ with the quality indicator:
\begin{equation}\label{eq:interaction}
  \ln Q_i \;=\; \alpha \;+\; \beta_0 \, \ln P_i \;+\; \beta_1 \, (\ln P_i \times \textit{Good}_i) \;+\; \gamma \, \textit{Good}_i \;+\; \mathbf{X}_i' \boldsymbol{\delta} \;+\; \mu_s \;+\; \varepsilon_i.
\end{equation}
If seedless and seeded marijuana occupy distinct market segments with different demand structures, the interaction coefficient $\beta_1$ will be nonzero.

\paragraph{Identification.}
The key identification challenge is that price may be correlated with unobserved demand-side factors.  In illicit markets, experienced users who purchase frequently may negotiate lower prices, generating a spurious positive correlation between quantity and the inverse of price that biases the OLS elasticity toward zero.  Conversely, buyers in high-demand locations may face elevated prices if local supply is inelastic, biasing the elasticity away from zero.

A second concern is measurement error in self-reported quantities and transaction values, especially because unit price is constructed as transaction value divided by quantity.  Asking value and quantity as separate transaction fields reduces reliance on a single reported unit-price measure, but it does not eliminate mechanical-ratio concerns if quantity is mismeasured.  We therefore report robustness checks that trim extreme price and quantity observations and exclude round-number quantities, the most visible form of heaping in the raw data.

These issues motivate an identification strategy that combines measurement-error checks with approaches to price endogeneity.  We first examine whether the estimated elasticity remains stable across progressively demanding fixed-effects structures, from no controls through state-by-month interactions.  We then implement three instrumental-variables strategies that exploit plausibly supply-driven variation in prices:

\begin{enumerate}
  \item \textbf{OSM driving time from production hubs.}  Following the supply-cost logic of \citet{davis2016crowd} and \citet{halcoussis2017geo}, we instrument price with the log driving time from the nearest major marijuana-producing state (Sinaloa, Chihuahua, Guerrero, Durango, and Michoac\'{a}n) to the purchase municipality.  The selection combines the traditional ``Golden Triangle'' production region of Sinaloa--Durango--Chihuahua with Guerrero and Michoac\'{a}n, which enforcement and cultivation assessments identify among Mexico's main marijuana-growing areas \citep{medel2015cultivation,ndic2008ndta}.  We calculate municipality centroids from projected INEGI municipal polygons and construct each production-hub point as the centroid of the corresponding projected full-state geometry.  We query OSRM/OpenStreetMap for source-to-market road-network driving times from those constructed hub points to purchase municipalities \citep{openstreetmap2026,luxen2011osrm}.  The relevant measure is therefore not a straight-line distance but the shortest-route driving time implied by the OSM road network:
  \[
    T_m = \min_{h \in \mathcal{H}} \text{time}_{h,m},
  \]
  where $m$ indexes the purchase municipality, $h$ indexes the five hub states, and $\text{time}_{h,m}$ is measured in driving minutes.  The direction is hub-to-market because the instrument is meant to proxy supply-chain transportation costs; in a directed road network, this need not equal the reverse trip.  Operationally, these are continuous origin--destination driving times: the pairwise analogue of isochrone travel-time bands, but without coarsening municipalities into broad polygons.  The instrument is $\ln T_m$.  The identifying assumption is that longer road-network driving times proxy for transportation and supply-chain costs that shift local retail prices, without directly affecting the quantity a given buyer demands conditional on controls and fixed effects.

  \item \textbf{Lagged seizure-gravity indexes.}  Using monthly municipality-level marijuana-seizure data compiled by \citet{mucd2019seizures} from official federal responses to public-information requests, we construct gravity-weighted supply-disruption measures:
  \[
    Z_{m,t-1} = \sum_j \frac{\text{kg}_{j,t-1}}{1+\tau_{j,m}},
  \]
  where $\text{kg}_{j,t-1}$ is the kilograms of marijuana seized in municipality $j$ during the month before the purchase, and $\tau_{j,m}$ is the OSRM/OpenStreetMap driving time in minutes from seizure municipality $j$ to purchase municipality $m$.  The one-month lag allows supply disruptions to propagate through local retail markets rather than assuming a same-month pass-through.  The $1+\tau_{j,m}$ denominator avoids division by zero and imposes a simple inverse-driving-time decay: seizures that are closer in the source-to-market road network receive more weight than seizures that are farther away in driving time.  Same-municipality seizures are retained; when $j=m$, the driving time is zero and the denominator equals one, so local seizures enter with full weight rather than being discarded.  We use $\ln(1+Z_{m,t-1})$ as the instrument.  The index is merged to each transaction by municipality and purchase month.  Because purchases in the analytic sample run from January through September~2019, the one-month lag uses MUCD marijuana-seizure records from December~2018 through August~2019 and can be merged to all 3{,}293 transactions.  Its identifying logic is that nearby or local interdiction activity can disrupt supply chains and shift retail prices, while being plausibly unrelated to the idiosyncratic quantity chosen by an individual buyer after conditioning on controls and fixed effects.

  We report two variants.  The hub-only version restricts $j$ to seizure municipalities in the five production-hub states used in the OSM driving-time instrument, so it asks whether enforcement shocks in upstream production areas predict retail prices across destination markets.  The stricter outside-region version excludes seizure municipalities in the buyer's own macroregion, reducing exposure to local demand shocks at the cost of weaker first-stage power.  Both variants use only marijuana seizures, not crop destruction or seed confiscations.

  \item \textbf{External Hausman-type instruments.}  Following the logic of Hausman-style price instruments used in differentiated-product demand settings \citep{hausman1994competitive,nevo2001cereal}, we compute mean log prices in other markets within the same quality--month cell.  The preferred version excludes all transactions from the buyer's own state; a stricter version excludes the buyer's entire macroregion.  This construction avoids using prices from the same local state market as the endogenous regressor, while preserving common supply-cost variation across markets for the same quality tier and month.  The corresponding IV columns also include quality-by-month fixed effects, so identification comes from cross-market price covariation net of national quality-specific monthly shocks.
\end{enumerate}

For each IV strategy, we report first-stage F-statistics and assess instrument strength.  We complement the IV analysis with Anderson--Rubin confidence sets that are robust to weak instruments.  Wu--Hausman p-values are reported only when the first stage clears the conventional $F \geq 10$ threshold, because the test is not informative under weak identification.

\paragraph{Robustness.}
We subject the baseline estimates to several additional checks, keeping the state-by-month fixed effects and demographic controls wherever the specification permits direct comparability: (i)~trimming at the 1st and 99th percentiles of price and quantity to reduce the influence of outliers; (ii)~excluding round-number quantities to address heaping-induced measurement error; (iii)~adding municipality fixed effects to absorb time-invariant local heterogeneity; (iv)~using wild cluster bootstrap inference to account for the limited number of state clusters; and (v)~the \citet{oster2019unobservable} coefficient-stability test, which quantifies the degree of selection on unobservables required to explain away the estimated elasticity.
